\documentclass[a4paper,12pt]{article}
\usepackage[utf8]{inputenc}
\usepackage[T1]{fontenc}
\usepackage[norsk]{babel}
\usepackage{amsmath, amssymb}
\usepackage{graphicx}
\usepackage{caption}
\usepackage{siunitx}
\usepackage{geometry}
\usepackage{float}
\usepackage{fancyhdr}
\usepackage{hyperref}

\geometry{margin=2.5cm}
\sisetup{output-decimal-marker={,}}

\title{Fysikkekspriment: \textit{Klinkekula}}
\author{Navn: André Hansen \\
Skole: ASK privatist \\
Data \today}
\date{}

\pagestyle{fancy}
\fancyhf{}
\rhead{Fysikkekspriment}
\lhead{André Hansen}
\cfoot{\thepage}

\begin{document}

\maketitle

\section*{Hensikt}

Hensikten med dette forsøket er å undersøke hvordan gravitasjon påvirker farten til en klinkekule når den sendes opp en skråfart. Jeg ønsker også å finne ut akelerasjonen og farten til klinkekulen.

Dette eksprimentet treffer flere kompetansemål innenfor fysikk:

\section*{Hypotese}

\end{document}



% \geometry{} % Removed or specify options, e.g. \geometry{margin=2cm}